\SourceLeaf{51}%
% A544-P0051-B02
X.\quad Si, manentibus reliquis ut in VII, $k$ est formae $8\mu+3$, atque $\frac{n}{a}$ formae $4\mu+3$, erit $\mathfrak A=-i$.

% A544-P0051-B03
XI.\quad Si, manentibus reliquis ut in VII, $k$ est formae $8\mu+7$, atque $\frac{n}{a}$ formae $4\mu+1$, erit $\mathfrak A=-i$.

% A544-P0051-B04
XII.\quad Si, manentibus reliquis ut in VII, $k$ est formae $8\mu+7$, atque $\frac{n}{a}$ formae $4\mu+3$, erit $\mathfrak A=+i$.

% A544-P0051-B05
Casum eum, ubi $a=2$, praeterimus; hic quidem $\mathfrak A$ foret $=\frac00$ sive indeterminatus, sed tunc semper $W=0$.

% A544-P0051-B06
Factores reliqui $\mathfrak B,\mathfrak C$ etc. perinde pendent a $b,c$ etc., ut $\mathfrak A$ ab $a$, quatenus in illorum determinationem ingrediuntur.

% A544-P0051-B07
\Article{31}

% A544-P0051-B08
Secundum hanc methodum alteram exemplum primum art. $29$ ita se habet:

% A544-P0051-B09
{\small
\[
\begin{aligned}
&\text{Factor }w\text{ fit }=(1+i)\surd2520\\
&\text{Pro }a=8\text{ factor respondens }\mathfrak A\text{ fit, per casum VIII, }=-1\\
&\text{Factori ipsius }n\text{ secundo }9\text{ respondet factor }+1\;\text{(per casum III.)}\\
&\text{Factori }5\text{ respondet factor }-1\;\text{(per casum II.)}\\
&\text{Factori }7\text{ respondet factor }-1\;\text{(per casum II.)}
\end{aligned}
\]
}

% A544-P0051-B10
Hinc conflatur productum $W=(-1-i)\surd2520$, ut in art. $29$.

% A544-P0051-B11
\Article{32}

% A544-P0051-B12
Quum valor ipsius $W$ per methodos \textit{duas} determinari possit, quarum altera relationibus numerorum $\frac{nk}{a}$, $\frac{nk}{b}$, $\frac{nk}{c}$ etc. ad numeros $a,b,c$ etc. innititur, altera vero a relationibus ipsius $k$ ad numeros $a,b,c$ etc. pendet, inter omnes has relationes nexus quidam conditionalis intercedere debet, ita ut quaevis e reliquis determinabilis esse debeat.\quad Supponamus, omnes numeros $a,b,c$ etc. esse numeros primos impares, atque $k$ accipi $=1$; distribuanturque factores $a,b,c$ etc. in duas classes, quarum altera contineat eos, qui sunt formae $4\mu+1$, et qui denotentur per $p,p',p''$ etc., altera vero constet ex iis, qui sunt formae $4\mu+3$, et qui exprimantur per $q,q',q''$ etc.: multitudinem posteriorum designabimus per $m$. His ita factis, observamus primo, $n$ fieri formae $4\mu+1$, si $m$ fuerit par (quorsum etiam referri debet casus is, ubi factores classis alterius omnino desunt, sive ubi $m=0$), contra $n$ fieri formae $4\mu+3$, si $m$ fuerit impar.\quad Iam determinatio
\SourceLeaf{52}%
% A544-P0052-B02
ipsius $W$ per methodum primam ita perficitur.\quad Pendeant numeri $P,P'\;P''$ etc., $Q,Q',Q''$ etc. ita a relationibus numerorum $\frac{n}{p}$, $\frac{n}{p'}$, $\frac{n}{p''}$ etc., $\frac{n}{q}$, $\frac{n}{q'}$, $\frac{n}{q''}$ etc. ad numeros $p,p',p''$ etc., $q,q',q''$ etc. resp., ut statuatur

% A544-P0052-B03
\[
\begin{aligned}
P&=+1,\;\text{si }\frac{n}{p}\text{ est residuum quadraticum ipsius }p\\
P&=-1,\;\text{si }\frac{n}{p}\text{ est non-residuum quadraticum ipsius }p
\end{aligned}
\]

% A544-P0052-B04
et perinde de reliquis.\quad Tunc erit $W$ productum e factoribus $P\surd p$, $P'\surd p'$, $P''\surd p''$ etc. $iQ\surd q$, $iQ'\surd q'$, $iQ''\surd q''$ etc., adeoque

% A544-P0052-B05
\[
\begin{aligned}
W=PP'P''\ldots QQ'Q''\ldots i^m\surd n
\end{aligned}
\]

% A544-P0052-B06
Per methodum secundam, aut potius statim per praecepta art. $19$, erit

% A544-P0052-B07
\[
\begin{aligned}
W&=+\surd n,\;\text{si }n\text{ est formae }4\mu+1,\\&\qquad\text{vel quod eodem redit, si }m\text{ est par}\\
W&=+i\surd n,\;\text{si }n\text{ est formae }4\mu+3,\\&\qquad\text{vel si }m\text{ est impar.}
\end{aligned}
\]

% A544-P0052-B08
Utrumque casum simul complecti licet per formulam sequentem:

% A544-P0052-B09
\[
\begin{aligned}
W=i^{mm}\surd n
\end{aligned}
\]

% A544-P0052-B10
Hinc itaque colligitur

% A544-P0052-B11
\[
\begin{aligned}
PP'P''\ldots QQ'Q''\ldots=i^{mm-m}
\end{aligned}
\]

% A544-P0052-B12
Sed $i^{mm-m}$ fit $=1$, quoties $m$ est formae $4\mu$ vel $4\mu+1$, atque $=-1$, quoties $m$ est formae $4\mu+2$ vel $4\mu+3$, unde deducimus sequens elegantissimum

% A544-P0052-B13
\textsc{Theorema.}\quad \textit{Denotantibus $a,b,c$ etc. numeros primos impares positivos inae}\textit{quales, quorum productum statuitur $=n$, et inter quos $m$ sint formae $4\mu+3$, re}\textit{liqui formae $4\mu+1$: multitudo eorum ex his numeris $a,b,c$ etc., quorum non-resi}\textit{dua resp. sunt $\frac{n}{a}$, $\frac{n}{b}$, $\frac{n}{c}$ etc., par erit, quoties $m$ est formae $4\mu$ vel $4\mu+1$, impar} \textit{vero, quoties $m$ est formae $4\mu+2$ vel $4\mu+3$.}

% A544-P0052-B14
Ita e. g. statuendo $a=3$, $b=5$, $c=7$, $d=11$, habemus tres numeros formae $4\mu+3$, puta $3$, $7$ et $11$; est autem $5\pdot7\pdot11\,R3$; $3\pdot7\pdot11\,R5$; $3\pdot5\pdot11\,R7$; $3\pdot5\pdot7\,N11$, sive unicus $\frac{n}{d}$ est non-residuum ipsius $d$.

% A544-P0052-B15
\Article{33}

% A544-P0052-B16
Celeberrimum \textit{theorema fundamentale} circa residua quadratica nihil aliud est, nisi casus specialis theorematis modo evoluti.\quad Limitando scilicet multitudinem
\SourceLeaf{53}%
% A544-P0053-B02
numerorum $a,b,c$ etc. ad \textit{duos}, patet, si unus tantum ex ipsis, vel neuter, sit formae $4\mu+3$, fieri debere vel simul $aRb$, $bRa$, vel simul $aNb$, $bNa$; contra si uterque est formae $4\mu+3$, unus ex ipsis alterius non-residuum esse debebit, atque hic illius residuum.\quad En itaque demonstrationem \textit{quartam} huius gravissimi theorematis, cuius demonstrationem primam et secundam in Disquisitionibus Arithmeticis, tertiam nuper in commentatione peculiari tradidimus (\textit{Commentt. T. XVI}): duas alias principiis rursus omnino diversis innitentes in posterum exponemus.\quad Summopere sane est mirandum, quod hocce venustissimum theorema, quod primo omnes conatus tam pertinaciter eluserat, tot postea viis toto coelo inter se distantibus adiri potuerit.

% A544-P0053-B03
\Article{34}

% A544-P0053-B04
Etiam theoremata reliqua, quae quasi supplementum ad theorema fundamentale efficiunt, scilicet per quae dignoscuntur numeri primi, quorum residua vel non-residua sunt $-1,+2$ et $-2$, ex iisdem principiis derivari possunt. Incipiemus a residuo $+2$.

% A544-P0053-B05
Statuendo $n=8a$, ita ut $a$ sit numerus primus, atque $k=1$, per methodum art. $28$. $W$ erit productum e duobus factoribus, quorum alter erit $+\surd a$, vel $+i\surd a$, si $8$, vel quod idem est $2$, est residuum quadraticum ipsius $a$; contra $-\surd a$ vel $-i\surd a$, si $2$ est non-residuum ipsius $a$.\quad Factor secundus autem est

% A544-P0053-B06
\[
\begin{aligned}
&(1+i)\surd8,\quad\text{si }a\text{ est formae }8\mu+1\\
&(-1+i)\surd8,\quad\text{si }a\text{ est formae }8\mu+3\\
&(-1-i)\surd8,\quad\text{si }a\text{ est formae }8\mu+5\\
&(1-i)\surd8,\quad\text{si }a\text{ est formae }8\mu+7
\end{aligned}
\]

% A544-P0053-B07
Sed per art. $18$ semper erit $W=(1+i)\surd n$; dividendo hunc valorem per quatuor valores factoris secundi, patet, factorem primum fieri debere

% A544-P0053-B08
\[
\begin{aligned}
&+\surd a,\quad\text{si }a\text{ est formae }8\mu+1\\
&-i\surd a,\quad\text{si }a\text{ est formae }8\mu+3\\
&-\surd a,\quad\text{si }a\text{ est formae }8\mu+5\\
&+i\surd a,\quad\text{si }a\text{ est formae }8\mu+7
\end{aligned}
\]

% A544-P0053-B09
Hinc sponte sequitur, in casu primo et quarto $2$ esse debere residuum ipsius $a$, in casu secundo et tertio autem non-residuum.

\SourceLeaf{54}%
% A544-P0054-B02
\Article{35}

% A544-P0054-B03
Numeri primi, quorum residuum vel non-residuum est $-1$, facile dignoscuntur adiumento theorematis sequentis, quod etiam per se ipsum satis memorabile est.

% A544-P0054-B04
\textsc{Theorema.}\quad \textit{Productum e duobus factoribus}

% A544-P0054-B05
\[
\begin{aligned}
W'&=1+r^{-1}+r^{-4}+\text{etc.}+r^{-(n-1)^2}\\
W&=1+r+r^4+\text{etc.}+r^{(n-1)^2}
\end{aligned}
\]

% A544-P0054-B06
\textit{est $=n$, si $n$ est impar; vel $=0$, si $n$ est impariter par; vel $=2n$, si $n$ est} \textit{pariter par.}

% A544-P0054-B07
\textit{Demonstr.}\quad Quum manifesto fiat

% A544-P0054-B08
\[
\begin{aligned}
W&=r+r^4+r^9+\text{etc.}+r^{nn}\\
&=r^4+r^9+\text{etc.}+r^{(n+1)^2}\\
&=r^9+\text{etc.}+r^{(n+2)^2}\quad\text{etc.}
\end{aligned}
\]

% A544-P0054-B09
productum $WW'$ ita quoque exhiberi poterit

% A544-P0054-B10
\[
\begin{aligned}
&1+r+r^4+r^9+\text{etc.}+r^{(n-1)^2}\\
+{}&r^{-1}(r+r^4+r^9+r^{16}+\text{etc.}+r^{nn})\\
+{}&r^{-4}(r^4+r^9+r^{16}+r^{25}+\text{etc.}+r^{(n+1)^2})\\
+{}&r^{-9}(r^9+r^{16}+r^{25}+r^{36}+\text{etc.}+r^{(n+2)^2})\\
&\text{etc.}\\
+{}&r^{-(n-1)^2}(r^{(n-1)^2}+r^{nn}+r^{(n+1)^2}\\&\qquad+r^{(n+2)^2}+\text{etc.}+r^{(2n-2)^2})
\end{aligned}
\]

% A544-P0054-B11
quod aggregatum verticaliter summatum producit

% A544-P0054-B12
\[
\begin{aligned}
&n\\
+{}&r(1+rr+r^4+r^6+\text{etc.}+r^{2n-2})\\
+{}&r^4(1+r^4+r^8+r^{12}+\text{etc.}+r^{4n-4})\\
+{}&r^9(1+r^6+r^{12}+r^{18}+\text{etc.}+r^{6n-6})\\
+{}&\text{etc.}\\
+{}&r^{(n-1)^2}(1+r^{2n-2}+r^{4n-4}\\&\qquad+r^{6n-6}+\text{etc.}+r^{2(n-1)^2})
\end{aligned}
\]

% A544-P0054-B13
Iam si $n$ impar est, singulae partes huius aggregati, praeter primam $n$, erunt $=0$; secunda enim manifesto fit $\frac{r(1-r^{2n})}{1-rr}$, tertia $\frac{r^4(1-r^{4n})}{1-r^4}$ etc. Quoties vero $n$ par est, excipere insuper oportebit partem

\SourceLeaf{55}%
% A544-P0055-B02
\[
\begin{aligned}
r^{\frac14nn}(1+r^n+r^{2n}+r^{3n}+\text{etc.}+r^{nn-n})
\end{aligned}
\]

% A544-P0055-B03
quae fit $=nr^{\frac14nn}$. In casu priori itaque fit $WW'=n$, in posteriori autem $=n+nr^{\frac14nn}$; sed $r^{\frac14nn}$ fit $=+1$, si $n$ est pariter par, tunc itaque prodit $WW'=2n$; contra fit $r^{\frac14nn}=-1$, si $n$ est impariter par, ubi itaque evadit $WW'=0$. Q. E. D.

% A544-P0055-B04
\Article{36}

% A544-P0055-B05
Iam per art. 22 constat, si $n$ sit numerus primus impar, $\frac{W'}{W}$ fieri $=+1$ vel $=-1$, prout $-1$ fuerit residuum vel non-residuum ipsius $n$. Hinc in casu priori esse debebit $W^2=+n$, in posteriori $W^2=-n$; quamobrem per art. 13 concludimus, casum priorem tunc tantum locum habere posse, quando $n$ sit formae $4\mu+1$, casumque posteriorem, quando $n$ sit formae $4\mu+3$.

% A544-P0055-B06
Denique e combinatione conditionum pro residuis $+2$ et $-1$ inventarum sponte sequitur, $-2$ esse residuum cuiusvis numeri primi formae $8\mu+1$ vel $8\mu+3$, atque non-residuum cuiusvis numeri primi formae $8\mu+5$ vel $8\mu+7$.

% A544-P0055-B07
\BookRule{42mm}

